% =====================================================================
% Rozdział 3: Reprezentacje i mapowania modelu
% =====================================================================

\chapter{Reprezentacje i mapowania modelu}
\label{chap:reprezentacje}

\section{Zasady reprezentacji modelu}
\label{sec:zasady-reprezentacji}
Model LEM definiuje znaczenie informacji o lokalizacji niezależnie od sposobu jej zapisu lub przekazu. Reprezentacja modelu określa sposób wyrażenia komponentów modelu LEM w określonym formacie danych lub mechanizmie wymiany informacji.

Wszystkie reprezentacje modelu LEM podlegają następującym zasadom:

\begin{itemize}
    \item \textbf{Niezmienność semantyki:} Reprezentacja musi wiernie odwzorowywać znaczenie komponentów zdefiniowanych w Rozdziale 2. Reprezentacja nie może wprowadzać nowych elementów semantycznych ani zmieniać definicji istniejących komponentów.
    \item \textbf{Jednoznaczność odwzorowania:} Każda reprezentacja musi umożliwiać jednoznaczne odwzorowanie zapisanych danych na komponenty modelu LEM.
    \item \textbf{Spójność z interpretacją kanoniczną:} Każda reprezentacja informacji LEM musi prowadzić do interpretacji zgodnej z interpretacją kanoniczną określoną w sekcji 2.4.5. Reprezentacja sprzeczna z interpretacją kanoniczną nie jest poprawną reprezentacją modelu LEM.
    \item \textbf{Równoważność reprezentacji:} Różne reprezentacje są semantycznie równoważne, jeżeli posiadają tę samą interpretację semantyczną zgodnie z zasadami określonymi w sekcji 2.4.3.
    \item \textbf{Autonomia formatu:} Reprezentacja może zawierać informacje dodatkowe, w tym metadane techniczne, o ile informacje te nie wpływają na interpretację komponentów modelu LEM.
    \item \textbf{Brak rozszerzenia modelu:} Reprezentacja lub mapowanie nie może rozszerzać znaczenia modelu LEM ani wprowadzać nowych komponentów informacji o lokalizacji. Dotyczy to również ograniczeń technicznych narzucanych przez format lub system zewnętrzny --- ograniczenie techniczne nigdy nie stanowi rozszerzenia semantyki modelu (zob. 4.5).
\end{itemize}

\begin{figure}[htbp]
\centering
\begin{tikzpicture}[node distance=9mm and 10mm,
  box/.style={draw,rounded corners,align=center,minimum width=23mm,minimum height=8mm},
  arr/.style={-{Latex[length=2mm]},thick}]
  \node[box] (core) {Rdzeń LEM};
  \node[box,above left=of core] (json) {JSON};
  \node[box,above=of core] (tlv) {TLV};
  \node[box,above right=of core] (yang) {YANG};
  \node[box,below left=of core] (civic) {Civic /\\PIDF-LO};
  \node[box,below=of core] (geo) {GeoJSON};
  \node[box,below right=of core] (dns) {DNS-safe\\alias};
  \draw[arr] (core) -- (json);
  \draw[arr] (core) -- (tlv);
  \draw[arr] (core) -- (yang);
  \draw[arr] (core) -- (civic);
  \draw[arr] (core) -- (geo);
  \draw[arr,dashed] (core) -- (dns);
\end{tikzpicture}
\caption{Rdzeń semantyczny i jego reprezentacje oraz mapowania}
\label{fig:lem-representations}
\end{figure}

\section{Reprezentacja tekstowa (Human-readable)}
\label{sec:reprezentacja-tekstowa}
Reprezentacja tekstowa służy do czytelnego przedstawienia informacji o lokalizacji zdefiniowanej przez model LEM.

Reprezentacja tekstowa przedstawia komponenty modelu LEM w formie przeznaczonej do interpretacji przez człowieka. Sposób rozmieszczenia informacji, kolejność występowania komponentów, stosowane separatory lub formatowanie graficzne nie wpływają na znaczenie opisu informacji o lokalizacji.

Znaczenie reprezentacji tekstowej wynika wyłącznie z obecnych komponentów modelu LEM oraz ich interpretacji semantycznej.

\noindent \textbf{Przykład informacyjny:}
\begin{verbatim}
Site: Main Campus
Building: Engineering Center
Level: 2
Reference Space: Room 201
Relative Position: East side
\end{verbatim}

Powyższy zapis przedstawia informację o lokalizacji w sposób czytelny dla człowieka. Konkretna forma prezentacji nie stanowi części modelu semantycznego LEM.

\section{Reprezentacja binarna (Compact / TLV)}
\label{sec:reprezentacja-binarna}
Reprezentacja binarna umożliwia wyrażenie informacji LEM w postaci zwartej struktury danych przeznaczonej do wymiany lub przechowywania informacji.

Mechanizm kodowania binarnego nie definiuje znaczenia informacji o lokalizacji. Znaczenie każdego elementu reprezentacji wynika z odwzorowania na komponenty modelu LEM.

Reprezentacja binarna może wykorzystywać mechanizmy takie jak Type-Length-Value (TLV) --- wzorzec kodowania rozpoznawalny formalnie m.in. z rekomendacji ITU-T X.690 \cite{ITU-X690} (zob. Załącznik B) --- jednak sposób przypisania typów, kodowania wartości oraz organizacji struktury danych należy do określonej specyfikacji reprezentacji lub profilu implementacyjnego.

W reprezentacji binarnej:
\begin{itemize}
    \item typ elementu reprezentacji odpowiada komponentowi modelu LEM,
    \item wartość elementu reprezentacji odpowiada wartości określonego komponentu,
    \item sposób kodowania wartości nie zmienia znaczenia komponentu.
\end{itemize}

Poprawna reprezentacja binarna musi prowadzić do tej samej interpretacji semantycznej, jaka wynikaby z dowolnej innej poprawnej reprezentacji informacji LEM.

\noindent \textbf{Przykład informacyjny} (profil kodowania TLV zdefiniowany poniżej: T = 1 oktet, L = 2 oktety big-endian, V = wartość; \textit{Level} używa jednego oktetu \texttt{kind} oraz czterech oktetów \texttt{int32} dla wartości numerycznej):

\begin{verbatim}
Typ (T)  Nazwa komponentu     Długość (L)  Wartość (V)
0x01     Site                 0x000B       "Main Campus"
0x02     Building             0x0012       "Engineering Center"
0x03     Level                0x0005       0x01 00000002
0x04     Reference Space      0x0007       "Room201"
0x05     Relative Position    0x0009       "East side"
\end{verbatim}

\noindent \textbf{Ciąg oktetów (zapis szesnastkowy):}
\begin{verbatim}
01 00 0B 4D 61 69 6E 20 43 61 6D 70 75 73
02 00 12 45 6E 67 69 6E 65 65 72 69 6E 67 20 43 65 6E 74 65 72
03 00 05 01 00 00 00 02
04 00 07 52 6F 6F 6D 32 30 31
05 00 09 45 61 73 74 20 73 69 64 65
\end{verbatim}

Powyższe przypisanie wartości typu (T) do konkretnego komponentu (np. \texttt{0x03} $\rightarrow$ Level) jest elementem konkretnej specyfikacji reprezentacji lub profilu implementacyjnego, a nie częścią modelu semantycznego LEM --- inny profil binarny może przypisać te same komponenty do innych wartości T, pozostając równoważny semantycznie (zob. 2.4.3), o ile odwzorowanie na komponenty modelu jest jednoznaczne.

\subsection{Formalny rejestr typów (informacyjny)}
\label{sec:rejestr-typow}
Poniższy rejestr formalizuje powyższy przykład:
\begin{enumerate}
    \item Pole L ma 2 oktety, co pozwala opisać pełny zakres długości reprezentacji UTF-8.
    \item Kodowanie Level odzwierciedla strukturę kind+value i wspólną dziedzinę \texttt{int32}.
    \item Relative Position jest zwykłym ciągiem UTF-8, bez podtypu.
\end{enumerate}

\begin{table}[h]
\centering
\small
\begin{tabular}{|l|l|p{8cm}|}
\hline
\textbf{T (1 oktet)} & \textbf{Komponent} & \textbf{Kodowanie wartości (V)} \\ \hline
0x01 & Site & UTF-8 \cite{RFC3629}, 1--1024 oktetów; po dekodowaniu 1--256 punktów kodowych Unicode \\ \hline
0x02 & Building & UTF-8 \cite{RFC3629}, 1--1024 oktetów; po dekodowaniu 1--256 punktów kodowych Unicode \\ \hline
0x03 & Level & 1 oktet kind (0x01=numeric, 0x02=non-building) + dla numeric dokładnie 4 oktety liczby \texttt{int32} ze znakiem, uzupełnienie do dwóch, big-endian; dla non-building brak dalszych oktetów \\ \hline
0x04 & Reference Space & UTF-8 \cite{RFC3629}, 1--1024 oktetów; po dekodowaniu 1--256 punktów kodowych Unicode \\ \hline
0x05 & Relative Position & UTF-8 \cite{RFC3629}, 1--1024 oktetów; po dekodowaniu 1--256 punktów kodowych Unicode \\ \hline
\end{tabular}
\end{table}

\noindent \textbf{Format nagłówka:} T (1 oktet) + L (2 oktety, big-endian, długość V w oktetach) + V.\\
Reguły spójności 2.5.1–2.5.4 stosuje się do zbioru zdekodowanych komponentów dokładnie tak samo, jak w przypadku reprezentacji JSON (zob. 2.5.6).

\section{Reprezentacja JSON}
\label{sec:reprezentacja-json}
Reprezentacja JSON umożliwia wyrażenie informacji LEM przy użyciu struktury danych JSON. JSON stanowi wyłącznie sposób reprezentacji komponentów modelu LEM i nie definiuje własnej semantyki informacji o lokalizacji.

Reprezentacja używa JSON zgodnego z \cite{RFC8259}. Parser MUSI odrzucać obiekty zawierające zduplikowane nazwy właściwości rdzenia, zanim zostanie utworzona postać kanoniczna.

Struktura JSON musi umożliwiać jednoznaczne odwzorowanie elementów reprezentacji na komponenty modelu LEM.

Nazwy właściwości, struktura obiektu oraz sposób kodowania wartości są elementami reprezentacji i nie mogą zmieniać znaczenia komponentów zdefiniowanych w modelu LEM.

Dwa dokumenty JSON są semantycznie równoważne, jeżeli ich interpretacja prowadzi do tej samej interpretacji kanonicznej modelu LEM.

Różnice w:
\begin{itemize}[noitemsep,topsep=2pt]
    \item kolejności właściwości obiektu JSON,
    \item sposobie formatowania dokumentu,
    \item technicznej reprezentacji wartości, jeżeli zachowana jest równoważność znaczenia,
\end{itemize}
nie wpływają na równoważność semantyczną informacji.

\noindent \textbf{Przykład informacyjny:}
\begin{verbatim}
{
  "site": "Main Campus",
  "building": "Engineering Center",
  "level": { "kind": "numeric", "value": 2 },
  "referenceSpace": "Room201",
  "relativePosition": "East side"
}
\end{verbatim}

Poniższy dokument jest semantycznie równoważny powyższemu (inna kolejność właściwości, inne formatowanie):
\begin{verbatim}
{
  "level": { "value": 2, "kind": "numeric" },
  "site": "Main Campus",
  "relativePosition": "East side",
  "referenceSpace": "Room201",
  "building": "Engineering Center"
}
\end{verbatim}

Nazwy właściwości użyte powyżej (\texttt{site}, \texttt{building}, \texttt{level} itd.) są elementem konkretnej reprezentacji --- profil lub schemat implementacyjny MOŻE stosować inne nazwy (zob. 3.1, 4.4.4), o ile odwzorowanie na komponenty modelu LEM pozostaje jednoznaczne.

\noindent \textbf{Przykład informacyjny --- przypadek non-building} (zob. 2.1.6): obiekt znajdujący się na dachu budynku:
\begin{verbatim}
{
  "site": "Main Campus",
  "building": "Engineering Center",
  "level": { "kind": "non-building", "value": null },
  "referenceSpace": "Roof"
}
\end{verbatim}

\section{Alias prezentacyjny DNS-safe}
\label{sec:reprezentacja-dns-safe}
Alias DNS-safe umożliwia prezentację wybranego zakresu informacji LEM w środowiskach posiadających ograniczenia charakterystyczne dla nazw DNS \cite{RFC1035}.

Ograniczenia techniczne wynikające z zasad tworzenia nazw DNS nie definiują znaczenia informacji o lokalizacji.

Znaczenie reprezentacji DNS-safe wynika wyłącznie z odwzorowania użytych elementów na komponenty modelu LEM.

Alias DNS-safe może przedstawiać jedynie podzbiór komponentów modelu LEM. Brak możliwości wyrażenia określonego komponentu w danym środowisku nie zmienia znaczenia modelu, lecz oznacza utratę informacji w aliasie.

Zdefiniowana niżej kanonikalizacja tekstu jest stratna, dlatego DNS-safe \textbf{NIE JEST} równoważną ani odwracalną reprezentacją LEM. Jest wyłącznie aliasem prezentacyjnym. System wymagający odtworzenia postaci kanonicznej MUSI przenosić pełną reprezentację LEM innym kanałem lub stosować oddzielne, bezstratne kodowanie z jednoznacznym schematem dekodowania.

\noindent \textbf{Przykład informacyjny:}
\begin{verbatim}
level2.room201.eng-center.main-campus.loc.example.
\end{verbatim}

\subsection{Formalny schemat etykiet (informacyjny)}
\label{sec:schemat-etykiet}
Formalny schemat wprowadza jawne, krótkie prefiksy w każdej etykiecie, niezależne od jej pozycji w nazwie:

\begin{table}[h]
\centering
\small
\begin{tabular}{|l|l|p{7cm}|}
\hline
\textbf{Prefiks} & \textbf{Komponent} & \textbf{Kodowanie rdzenia etykiety} \\ \hline
s- & Site & Kanonikalizowany ciąg znaków \\ \hline
b- & Building & Kanonikalizowany ciąg znaków \\ \hline
l- & Level & \texttt{l-<N>} dla kind="numeric", wartość nieujemna; \texttt{l-n<N>} dla wartości ujemnej; \texttt{l-nb} dla kind="non-building" \\ \hline
r- & Reference Space & Kanonikalizowany ciąg znaków \\ \hline
\end{tabular}
\end{table}

\noindent Relative Position nigdy nie jest reprezentowany w DNS-safe.\\
Struktura pełnej nazwy (notacja ABNF, zgodnie z \cite{RFC5234}):
\begin{verbatim}
lem-dns-name  = *(component-label ".") suffix "."
component-label = prefix "-" label-core
prefix        = "s" / "b" / "l" / "r"
label-core    = 1*61(ALPHA / DIGIT / "-")
suffix        = 1*(dns-label ".") dns-label
\end{verbatim}

Kanonikalizacja wartości tekstowej (Site, Building, Reference Space) do rdzenia etykiety: zamiana na małe litery, zastąpienie spacji i innych znaków spoza \texttt{[a-z0-9]} znakiem \texttt{-}, zwinięcie kolejnych \texttt{-} do jednego, usunięcie \texttt{-} z brzegów. Ta kanonikalizacja jest stratna.

\noindent \textbf{Uwaga:} Punycode \cite{RFC3492} koduje Unicode w etykietach DNS, ale samo jego użycie nie odtwarza informacji utraconych przez usuwanie separatorów, zakres komponentów ani pominięcie \textit{Relative Position}. Bezstratny wariant wymaga osobnej specyfikacji.

\noindent \textbf{Przykład formalny:}
\begin{verbatim}
s-main-campus.b-engineering-center.l-2.r-room201.loc.example.
\end{verbatim}

\noindent \textbf{Przykład z Level ujemnym:}
\begin{verbatim}
s-main-campus.l-n1.r-parking-p1.loc.example.
\end{verbatim}

\noindent \textbf{Przykład z non-building (dach):}
\begin{verbatim}
s-main-campus.b-engineering-center.l-nb.r-roof.loc.example.
\end{verbatim}

\section{Mapowanie do istniejących modeli informacji}
\label{sec:mapowanie-modeli}
Istniejące modele informacji mogą być wykorzystywane do wyrażania wybranych aspektów informacji LEM poprzez zdefiniowane mapowania semantyczne.

Mapowanie do zewnętrznego modelu informacji nie oznacza zastąpienia modelu LEM ani rozszerzenia jego znaczenia.

Zewnętrzny model pełni rolę źródła reprezentacji określonych komponentów, których znaczenie musi być zgodne z interpretacją LEM.

Komponenty zewnętrznego modelu, które nie posiadają odpowiednika w modelu LEM, pozostają informacją dodatkową i nie stanowią części opisu informacji LEM.

\subsection{Mapowanie do modeli YANG}
\label{sec:mapowanie-yang}
Modele YANG mogą być wykorzystywane jako reprezentacja strukturalna komponentów LEM w środowiskach zarządzania sieciowego i systemowego.

\noindent \textbf{Przykład informacyjny --- poglądowy moduł YANG \cite{RFC7950}:}
\begin{verbatim}
module example-lem-location {
  namespace "urn:example:lem-location";
  prefix "lem";
  description
    "Poglądowe odwzorowanie komponentów modelu semantycznego LEM (Rozdział 2)
     na węzły danych YANG. Wyłącznie informacyjne (zob. 3.6.1).
     Struktura Level odzwierciedla 2.1.6 (kind/value).";

  grouping lem-location {
    description "Grupa węzłów odpowiadająca komponentom modelu LEM (zob. 2.1).";

    leaf site {
      type string { length "1..256"; }
      description "Odwzorowanie komponentu LEM: Site (zob. 2.1.6).";
    }

    leaf building {
      type string { length "1..256"; }
      description "Odwzorowanie komponentu LEM: Building (zob. 2.1.6).";
    }

    container level {
      description
        "Odwzorowanie komponentu LEM: Level (zob. 2.1.6). Struktura
         kind/value zamiast prostej liczby całkowitej.";

      leaf kind {
        type enumeration {
          enum numeric { description "Kondygnacja numerowana."; }
          enum non-building { description "Dach, parking naziemny, przestrzeń zewnętrzna itp."; }
        }
        mandatory true;
      }

      leaf value {
        type int32;
        when "../kind = 'numeric'";
        description "Wymagane i całkowite dla kind=numeric; nieobecne dla kind=non-building.";
      }

      must "(kind = 'numeric' and value) or (kind = 'non-building' and not(value))" {
        error-message "Zgodnie z 2.1.6: value wymagane dla kind=numeric, niedozwolone dla kind=non-building.";
      }
    }

    leaf reference-space {
      type string { length "1..256"; }
      description "Odwzorowanie komponentu LEM: Reference Space (zob. 2.1.6).";
    }

    leaf relative-position {
      type string { length "1..256"; }
      description "Odwzorowanie komponentu LEM: Relative Position (zob. 2.1.6) — opis relacyjny.";
      must "../reference-space" {
        error-message "Relative Position obecny bez Reference Space narusza LEM 2.5.3.";
      }
    }
  }

  container device-location {
    description "Przykładowy węzeł konfiguracyjny wykorzystujący grupę lem-location.";
    uses lem-location;
  }
}
\end{verbatim}

\subsection{Mapowanie do modeli Civic Address}
\label{sec:mapowanie-civic}
Modele Civic Address mogą być wykorzystywane do reprezentacji wybranych komponentów opisujących adresowe aspekty informacji o lokalizacji.

\noindent \textbf{Przykład informacyjny --- dokument PIDF-LO civicAddress \cite{RFC5139}:}
\begin{verbatim}
<civicAddress xml:lang="en-US"
 xmlns="urn:ietf:params:xml:ns:pidf:geopriv10:civicAddr">
  <country>US</country>
  <NAM>Main Campus</NAM>
  <BLD>Engineering Center</BLD>
  <FLR>2</FLR>
  <ROOM>Room201</ROOM>
  <LOC>East side</LOC>
</civicAddress>
\end{verbatim}

\begin{table}[h]
\centering
\small
\begin{tabularx}{\textwidth}{|>{\raggedright\arraybackslash}p{2.2cm}|>{\raggedright\arraybackslash}p{3.3cm}|>{\raggedright\arraybackslash}X|}
\hline
\textbf{CAtype} & \textbf{Komponent LEM} & \textbf{Charakter mapowania} \\ \hline
BLD (27) & Building & Odwzorowanie bezpośrednie. \\ \hline
FLR (28) & Level, kind="numeric" & Odwzorowanie dotyczy wyłącznie przypadku kind="numeric". \\ \hline
(brak wprost) & Level, kind="non-building" & Civic Address nie posiada elementu dedykowanego. \\ \hline
ROOM (23) & Reference Space & Odwzorowanie bezpośrednie. \\ \hline
NAM (22) & Site (interpretacja) & Interpretacyjne. \\ \hline
LOC (1) & Relative Position & Semantycznie bliskie, lecz nie identyczne. \\ \hline
\end{tabularx}
\caption{Mapowanie wybranych typów Civic Address \cite{RFC4776} na komponenty LEM}
\end{table}

\subsection{Mapowanie do modeli geospatial}
\label{sec:mapowanie-geospatial}
Modele geospatial mogą być wykorzystywane do wyrażania przestrzennych aspektów informacji o lokalizacji, w szczególności współrzędnych lub odniesień przestrzennych.

\noindent \textbf{Przykład informacyjny --- obiekt GeoJSON \cite{RFC7946}:}
\begin{verbatim}
{
  "type": "Feature",
  "geometry": {
    "type": "Point",
    "coordinates": [-73.961452, 40.808058]
  },
  "properties": {
    "lem:site": "Main Campus",
    "lem:building": "Engineering Center",
    "lem:level": { "kind": "numeric", "value": 2 },
    "lem:referenceSpace": "Room201",
    "lem:relativePosition": "East side"
  }
}
\end{verbatim}

\noindent Prefiks \texttt{lem:} w nazwach właściwości jest konwencją informacyjną. Kolejność współrzędnych \texttt{[longitude, latitude]} jest wymagana przez \cite{RFC7946} §4.

\section{Przykład ilustracyjny (jedna lokalizacja we wszystkich reprezentacjach)}
\label{sec:przyklad-ilustracyjny}
Poniższy przykład przedstawia tę samą informację o lokalizacji wyrażoną kolejno we wszystkich reprezentacjach. Wszystkie poniższe zapisy MUSZĄ prowadzić do tej samej interpretacji kanonicznej (zob. 2.4.5).

\noindent \textbf{Opisywana lokalizacja:} obiekt znajdujący się w budynku "Engineering Center" na kampusie "Main Campus", na drugiej kondygnacji, w pomieszczeniu "Room201", po wschodniej stronie tego pomieszczenia.

\begin{enumerate}
    \item \textbf{Reprezentacja tekstowa (3.2):}
\begin{verbatim}
Site: Main Campus
Building: Engineering Center
Level: 2 (numeric)
Reference Space: Room201
Relative Position: East side
\end{verbatim}

    \item \textbf{Reprezentacja binarna / TLV (3.3, wg formalnego rejestru 3.3.1 — pole L na 2 oktetach):}
\begin{verbatim}
01 00 0B 4D 61 69 6E 20 43 61 6D 70 75 73
02 00 12 45 6E 67 69 6E 65 65 72 69 6E 67 20 43 65 6E 74 65 72
03 00 05 01 00 00 00 02
04 00 07 52 6F 6F 6D 32 30 31
05 00 09 45 61 73 74 20 73 69 64 65
\end{verbatim}

    \item \textbf{Reprezentacja JSON (3.4):}
\begin{verbatim}
{
  "site": "Main Campus",
  "building": "Engineering Center",
  "level": { "kind": "numeric", "value": 2 },
  "referenceSpace": "Room201",
  "relativePosition": "East side"
}
\end{verbatim}

    \item \textbf{Alias prezentacyjny DNS-safe (3.5.1) --- zapis stratny:}
\begin{verbatim}
s-main-campus.b-engineering-center.l-2.r-room201.loc.example.
\end{verbatim}
\end{enumerate}

\section{Przykłady niepoprawnych opisów lokalizacji}
\label{sec:przyklady-niepoprawne}
\begin{itemize}
    \item \textbf{Przykład 1 --- naruszenie spójności relacji przestrzennych (2.5.3):}
\begin{verbatim}
{
  "site": "Main Campus",
  "building": "Engineering Center",
  "level": { "kind": "numeric", "value": 2 },
  "relativePosition": "East side"
}
\end{verbatim}
Obecny jest komponent \textit{Relative Position}, lecz brakuje \textit{Reference Space}.

    \item \textbf{Przykład 2 --- pusty opis lokalizacji (2.4.4):}
\begin{verbatim}
{}
\end{verbatim}

    \item \textbf{Przykład 3 --- naruszenie struktury Level (2.1.6):}
\begin{verbatim}
{
  "site": "Main Campus",
  "level": { "kind": "non-building", "value": 3 }
}
\end{verbatim}
\end{itemize}
